\section{Teori Informasi Shannon}

\subsection{Entropi sebagai Ukuran Ketidakpastian}

Claude Shannon dalam paper "A Mathematical Theory of Communication" (1948) memperkenalkan konsep \textbf{entropi} sebagai ukuran matematis untuk ketidakpastian atau informasi \cite{shannon1948}. Entropi menjadi fondasi teoretis untuk memahami keamanan kriptografi.

\textbf{Definisi Entropi}: Untuk sumber informasi $X$ dengan alfabet $\mathcal{X}$ dan distribusi probabilitas $p(x)$, entropi $H(X)$ didefinisikan sebagai:
\[
H(X) = -\sum_{x \in \mathcal{X}} p(x) \log_2 p(x) \text{ bit}
\]

Entropi mengukur rata-rata informasi per simbol yang dihasilkan oleh sumber. Semakin tinggi entropi, semakin tidak terprediksi sumber tersebut.

\begin{contoh}
Untuk alfabet Inggris 26 huruf dengan distribusi seragam:
\[
H(X) = -\sum_{i=1}^{26} \frac{1}{26} \log_2 \frac{1}{26} = \log_2 26 \approx 4.7 \text{ bit/huruf}
\]

Untuk alfabet dengan distribusi tidak seragam (seperti bahasa Inggris sebenarnya), entropi lebih rendah karena beberapa huruf lebih sering muncul.
\end{contoh}

\subsection{Laju Bahasa dan Redundansi}

\textbf{Laju Bahasa} (language rate) adalah entropi per simbol dalam bahasa alami. Bahasa Inggris memiliki laju bahasa sekitar 1.5 bit/huruf, jauh lebih rendah dari entropi maksimum 4.7 bit/huruf \cite{shannon1951}.

\textbf{Redundansi} $R$ didefinisikan sebagai:
\[
R = 1 - \frac{H_{\text{actual}}}{H_{\text{maximum}}}
\]

Untuk bahasa Inggris:
\[
R = 1 - \frac{1.5}{4.7} \approx 0.68 \text{ (68\% redundansi)}
\]

Redundansi ini adalah alasan mengapa \textbf{frequency analysis} dapat memecahkan cipher substitusi sederhana.

\begin{table}[htbp]
  \centering
  \small
  \begin{tabular}{lcc}
    \toprule
    Sumber & Entropi Maksimum & Entropi Aktual \\
    \midrule
    Alfabet seragam (26 huruf) & 4.7 bit & 4.7 bit \\
    Bahasa Inggris & 4.7 bit & 1.5 bit \\
    Binary seragam & 1.0 bit & 1.0 bit \\
    \bottomrule
  \end{tabular}
  \caption{Perbandingan entropi maksimum vs aktual.}
\end{table}

\subsection{Unicity Distance}

\textbf{Unicity distance} adalah panjang ciphertext minimum di mana hanya ada satu plaintext yang masuk akal yang dapat menghasilkan ciphertext tersebut \cite{shannon1949}. Ini adalah ukuran penting untuk keamanan cipher.

Untuk cipher dengan kunci $K$ dan plaintext $P$, unicity distance $U$ didefinisikan sebagai:
\[
U = \frac{H(K)}{D}
\]
dimana $H(K)$ adalah entropi kunci dan $D$ adalah redundansi per simbol.

\begin{contoh}
Untuk Caesar cipher:
\begin{itemize}
  \item Ruang kunci: 25 kemungkinan
  \item $H(K) = \log_2 25 \approx 4.64$ bit
  \item Redundansi bahasa Inggris: $D = 4.7 - 1.5 = 3.2$ bit/huruf
  \item $U = \frac{4.64}{3.2} \approx 1.45$ huruf
\end{itemize}

Ini berarti hanya perlu 2 huruf ciphertext untuk memecahkan Caesar cipher secara unik!
\end{contoh}

\begin{table}[htbp]
  \centering
  \small
  \begin{tabular}{lcc}
    \toprule
    Cipher & Ukuran Kunci & Unicity Distance \\
    \midrule
    Caesar (25 kunci) & 4.64 bit & 2 huruf \\
    Vigenère (26$^k$ kunci) & $k \cdot 4.7$ bit & $1.5k$ huruf \\
    DES (56-bit kunci) & 56 bit & 17.5 huruf \\
    AES-128 (128-bit kunci) & 128 bit & 40 huruf \\
    \bottomrule
  \end{tabular}
  \caption{Unicity distance untuk berbagai cipher.}
\end{table}

\subsection{Implikasi untuk Keamanan Kriptografi}

Teori informasi Shannon memberikan wawasan fundamental tentang keamanan kriptografi:

\begin{enumerate}
  \item \textbf{Kerentanan Frequency Analysis}: Redundansi tinggi dalam bahasa alami membuat cipher substitusi rentan terhadap frequency analysis.
  
  \item \textbf{Kebutuhan Kunci yang Cukup}: Unicity distance menunjukkan bahwa cipher dengan kunci pendek tidak aman untuk pesan panjang.
  
  \item \textbf{Diffusion dan Confusion}: Shannon mengusulkan dua properti untuk cipher yang aman:
  \begin{itemize}
    \item \textbf{Confusion}: Menyembunyikan hubungan antara kunci dan ciphertext
    \item \textbf{Diffusion}: Menyebar pengaruh satu plaintext bit ke banyak ciphertext bits
  \end{itemize}
  
  \item \textbf{Perfect Secrecy}: One-Time Pad mencapai perfect secrecy karena $H(K) \geq H(P)$.
\end{enumerate}

\subsection{Perfect Secrecy dan One-Time Pad}

\textbf{Perfect secrecy} terjadi ketika ciphertext tidak memberikan informasi apapun tentang plaintext. Shannon membuktikan bahwa perfect secrecy hanya mungkin jika:
\[
H(K) \geq H(P)
\]

One-Time Pad memenuhi syarat ini karena:
\begin{itemize}
  \item Kunci acak dengan panjang sama dengan plaintext
  \item Kunci hanya digunakan sekali
  \item $H(K) = n$ bit untuk plaintext $n$ bit
  \item $H(P) \leq n$ bit
\end{itemize}

\begin{contoh}
Dengan plaintext 8-bit (ASCII 'A' = 65 = 01000001) dan kunci acak 8-bit:
\begin{itemize}
  \item Plaintext: 01000001
  \item Kunci: 11010010 (acak)
  \item Ciphertext: 10010011 (XOR)
\end{itemize}
Setiap kunci yang mungkin (256 kemungkinan) menghasilkan ciphertext yang berbeda, sehingga ciphertext tidak memberikan informasi tentang plaintext.
\end{contoh}

\subsection{Distance of Uniqueness dalam Praktik}

Unicity distance memiliki implikasi praktis penting:

\begin{itemize}
  \item \textbf{DES}: $U \approx 17.5$ huruf, sangat rentan terhadap known-plaintext attack
  \item \textbf{AES-128}: $U \approx 40$ huruf, lebih aman untuk pesan pendek
  \item \textbf{RSA 2048-bit}: $U \approx 640$ huruf, sangat aman untuk pesan praktis
\end{itemize}

\begin{catatan}
Unicity distance adalah batas teoretis minimum. Dalam praktik, serangan sering berhasil dengan ciphertext yang lebih pendek karena informasi tambahan (struktur bahasa, format pesan, dll) dapat membantu penyerang.
\end{catatan}

\subsection{Koneksi ke Bab-Bab Berikutnya}

Konsep teori informasi ini akan menjadi dasar untuk:
\begin{itemize}
  \item \textbf{Bab III}: Memahami mengapa matematika penting untuk keamanan
  \item \textbf{Bab IV}: Analisis kelemahan cipher klasik menggunakan frequency analysis
  \item \textbf{Bab V}: Desain cipher modern dengan diffusion dan confusion
  \item \textbf{Bab VI}: Pemilihan ukuran kunci yang aman berdasarkan unicity distance
\end{itemize}

Pemahaman teori informasi memberikan kerangka matematis untuk mengevaluasi keamanan cipher dan merancang sistem kriptografi yang lebih kuat.
